% Zone „Das kennst du schon“ – aus bank/wahrscheinlichkeit/zone.jsonl (zusammenbau v0.1)
\einheitenkopf[zone][Kennst du schon]{Das kennst du schon}
\setcounter{aufgabe}{0}

%% TODO Titel: Ich-kann-Satz fehlt in der Bank (Platzhalter: Kettenname) – Nr. 1
%% TODO Anweisung: ein Satz über der ersten Teilaufgabe fehlt in der Bank – Nr. 1
\begin{aufgabe}{Brüche kürzen; Bruch ↔ Dezimalzahl ↔ Prozent (Achtelbruch als Dezimalzahl und Prozentsatz)}
\begin{teile}
\teil Kürze $\frac{6}{8}$ und schreibe den Bruch als Prozentsatz.
\teil Schreibe $\frac{9}{15}$ gekürzt, als Dezimalzahl und als Prozentsatz.
\end{teile}
\end{aufgabe}

%% TODO Titel: Ich-kann-Satz fehlt in der Bank (Platzhalter: Kettenname) – Nr. 2
%% TODO Anweisung: ein Satz über der ersten Teilaufgabe fehlt in der Bank – Nr. 2
\begin{aufgabe}{Anteil einer Menge nehmen (30 \% von 10 Feldern; zwei Drittel von 9 Kugeln)}
\begin{teile}
\teil Wie viel ist $\frac{1}{3}$ von 9 Kugeln?
\teil Wie viel sind $\frac{2}{5}$ von 45 Losen?
\end{teile}
\end{aufgabe}

%% TODO Titel: Ich-kann-Satz fehlt in der Bank (Platzhalter: Kettenname) – Nr. 3
%% TODO Anweisung: ein Satz über der ersten Teilaufgabe fehlt in der Bank – Nr. 3
\begin{aufgabe}{Relative Häufigkeit berechnen (Anzahl geteilt durch Gesamtzahl)}
\begin{teile}
\teil Tim wirft 8-mal auf den Korb und trifft 4-mal. Wie groß ist die relative Häufigkeit der Treffer?
\teil Bei 80 Drehungen eines Glücksrads kommt 14-mal Rot. Wie groß ist die relative Häufigkeit für Rot in Prozent?
\end{teile}
\end{aufgabe}

%% TODO Titel: Ich-kann-Satz fehlt in der Bank (Platzhalter: Kettenname) – Nr. 4
%% TODO Anweisung: ein Satz über der ersten Teilaufgabe fehlt in der Bank – Nr. 4
\begin{aufgabe}{Anzahl der Zahlen in einem Bereich (101 bis 900 sind 800 Lose; Endziffer 6 in 106 bis 896 sind 80)}
\begin{teile}
\teil Lose tragen die Nummern 1 bis 60. Wie viele Lose sind es?
\teil Wie viele der Zahlen von 1 bis 90 enden auf 4?
\end{teile}
\end{aufgabe}

%% TODO Titel: Ich-kann-Satz fehlt in der Bank (Platzhalter: Kettenname) – Nr. 5
%% TODO Anweisung: ein Satz über der ersten Teilaufgabe fehlt in der Bank – Nr. 5
\begin{aufgabe}{Brüche multiplizieren (auch drei Faktoren) und gleichnamige Brüche addieren}
\begin{teile}
\teil Berechne $\frac{1}{2} \cdot \frac{1}{3}$.
\teil Berechne $\frac{2}{3} \cdot \frac{3}{4} \cdot \frac{1}{2}$ und kürze.
\end{teile}
\end{aufgabe}

%% TODO Titel: Ich-kann-Satz fehlt in der Bank (Platzhalter: Kettenname) – Nr. 6
%% TODO Anweisung: ein Satz über der ersten Teilaufgabe fehlt in der Bank – Nr. 6
\begin{aufgabe}{Anzahl der Zahlen in einem Bereich (101 bis 900 sind 800 Lose; Endziffer 6 in 106 bis 896 sind 80) – Fehler finden}
\begin{teile}
\teil Ben berechnet die Zahl der Lose mit den Nummern 301 bis 700 so: $700 - 301 = 399$. Finde den Fehler und rechne richtig.
\end{teile}
\end{aufgabe}

%% TODO Titel: Ich-kann-Satz fehlt in der Bank (Platzhalter: Kettenname) – Nr. 7
%% TODO Anweisung: ein Satz über der ersten Teilaufgabe fehlt in der Bank – Nr. 7
\begin{aufgabe}{Anzahl der Zahlen in einem Bereich (101 bis 900 sind 800 Lose; Endziffer 6 in 106 bis 896 sind 80) – selbst rechnen}
\begin{teile}
\teil Lose tragen die Nummern 401 bis 950. Wie viele Lose sind es?
\end{teile}
\end{aufgabe}
